\documentclass[11pt]{article}
\usepackage[margin=1in]{geometry}
\usepackage{amsmath,amssymb,mathtools}
\usepackage[unicode,colorlinks=true,linkcolor=blue,urlcolor=blue]{hyperref}
\usepackage{microtype}
\providecommand{\tightlist}{}
\setlength{\emergencystretch}{2em}
\title{Consecutive excedance zeros of permutations with a minimum cycle length}
\author{ProgramComputer}
\date{November 25, 2026}
\begin{document}
\maketitle
\begin{abstract}
We study the excedance enumerator for permutations whose cycles all have at least $k$ letters. After removing the common zero, consecutive rows strictly interlace at every size when $k=3$, and whenever $n\ge Bk^2$ for a single absolute constant $B$. A comparison of consecutive coefficients of the exponential generating function controls the middle zeros. A decreasing-run model controls the outer zeros. For $k=3$ the analytic estimates meet at size 400, and an integer certificate handles the finite prefix. A separated-pole expansion and an explicit bound on the numerator of the generating function prove strict interlacing for all $n\ge10^{20}$, uniformly for all $3\le k\le n$. The remaining finite range is too large for exhaustive verification.
We also establish strict interlacing in every size for which the next permutation can contain at most two cycles. Here the analytic estimates start at size 192 and an exact finite replay treats every smaller pair.
\end{abstract}
\textbf{Theorem 1.} Let $E_{n,k}(t)$ count permutations of $[n]$ with all cycles of length at least $k$, by their number of excedances. For $k=3$ and every $n\ge3$, the $n-2$ nonzero zeros of $E_{n,3}$ are simple and negative. They strictly interlace the nonzero zeros of $E_{n+1,3}$, with the next row supplying both outer roots. More generally, for every fixed $k\ge3$ the same assertions hold for all sufficiently large $n$. Theorem 2 gives these conclusions for all $n\ge Bk^2$ with one absolute constant $B$. In all these rows zero itself is simple. This treats the first conjectured threshold in Alexandersson’s statement and reduces every other fixed threshold to finitely many rows. The proof gives an explicit analytic cutoff of 400 at $k=3$. \href{https://www.symmetricfunctions.com/realRootedOpenProblems.htm\#conjLongCycleExcedances}{Original conjecture} · \href{https://github.com/PerAlexandersson/RealRooted/issues/2590\#issuecomment-7059794955}{Author’s derivation and computations}

\section{Poles and the two adjacent coefficients}

Write $t=-e^{-s}$ for $s\ge0$. The symmetry $E_{n,k}(t)=t^n E_{n,k}(1/t)$ follows by permutation inversion, since a derangement has no fixed points. Its coefficients of $t$ and $t^{n-1}$ are one. We first work on $-1\le t<0$ and then use this symmetry.

Let $A_m(t)$ be the ordinary Eulerian polynomial with $A_0=1$. A cycle of length $l\ge2$ has enumerator $t A_{l-1}(t)$, by writing the largest label first and counting ascents in the cycle. A singleton has weight one. The unrestricted permutation EGF is $(1-t)e^w/(1-te^w)$ in the variable $w=(1-t)z$. For instance this follows by expanding the denominator and using $\sum_{j\ge0}(j+1)^n t^j=A_n(t)/(1-t)^{n+1}$. By the exponential formula, deleting cycles of lengths below $k$ therefore gives

\[
 \frac{E_{n,k}(t)}{n!(1-t)^{n+1}}=[w^n]\frac{T_s(w)}{1+e^{w-s}},\qquad
 T_s(w)=\exp\!\left(a w-\sum_{l=2}^{k-1}
 \frac{t A_{l-1}(t)}{l!(1-t)^l}w^l\right),\quad a=\frac{-t}{1-t}.
\]

The Eulerian power-sum identity used here also follows by applying $1+t\,d/dt$ repeatedly to $(1-t)^{-1}$. At $k=3$, put $b=a(1-a)$. Then $T_s(w)=\exp(aw+bw^2/2)$.

The poles are $w_p=s+i p\pi$ for odd integers $p$. Let $r=|w_1|$, $\theta=\arg w_1$, $S=\sin\theta=\pi/r$ and $z_p=r/w_p=(\cos\theta+i pS)^{-1}$. Let $\gamma(\theta)$ be the imaginary part of the displayed exponent of $T_s(w_1)$. Define $\Psi_n=(n+1)\theta-\gamma(\theta)$ and

\[
 Y_n=\frac{r^{n+1} E_{n,k}(t)}{2|T_s(w_1)|n!(1-t)^{n+1}}.
\]

The pair $p=\pm1$ contributes $\cos\Psi_n$ to $Y_n$. More precisely, take

\[
 v=2\pi\max(1,\lceil s/(4\pi)\rceil),\qquad R=\sqrt{s^2+v^2}.
\]

Deforming the coefficient circle to $|w|=R$ adds $T_s(w_p)/w_p^{n+1}$ for each enclosed pole: the residue of the generating function there is $-T_s(w_p)$. In $Y_n$ this is the sum over positive odd $p$ with $p\pi<v$ of $\operatorname{Re}(T_s(w_p)z_p^{n+1})/|T_s(w_1)|$, plus a circle error $e_n$. The same contour works for $n+1$.

Here $v\ge\max(2\pi,s/2)$, so $(r/R)^2\le17/20$, with the worst bound at $s=4\pi$. Every point of the circle has distance at least $1/2$ from every pole. Indeed, the nearest radii are those with $|p|\pi=v\pm\pi$; their radial distance is at least $v\pi/(2R+\pi)\ge\pi/(2\sqrt5+1/2)>1/2$. The inverse denominator has modulus at most $8$ on the circle. If $|\Re(w)-s|\ge1/4$, use $|1-e^{-1/4}|>1/5$. Otherwise the distance of $\Im(w)$ to the nearest odd multiple of $\pi$ exceeds $2/5$, and $|1+e^{x+iy}|^2=(1-e^x)^2+4e^x\sin^2(\operatorname{dist}(y,(2\mathbb Z+1)\pi)/2)$ with $x=\Re(w)-s$ gives the bound by $\sin u\ge 2u/\pi$ on $[0,\pi/2]$.

For $k=3$, $|T_s(w_1)|\ge e^{-\pi^2/8}>1/4$. On $|w|\le R$ we have $R\le 3s/2+2\pi$, $a\le\min(1/2,e^{-s})$ and $b\le\min(1/4,e^{-s})$. Hence

\[
 \log |T_s(w)|\le aR+bR^2/2
 \le 3/(2e)+\pi+9/(2e^2)+3\pi/e+\pi^2/2<14 .
\]

For $n\ge400$ and $nS^2\ge1$ the circle estimate yields

\[
 |e_n|,|e_{n+1}|\le 16 e^{14} r(17/20)^{n/2}< S/100.
\]

Indeed $16 e^{14}\pi n(17/20)^{n/2}<1/100$ at $400$ and decreases. This inequality can also be checked with the rational upper bounds $e<3$ and $\pi<22/7$.

We bound the other enclosed poles. For $n\ge100$ and $nS^2\ge1$ the monotonicity of $(1+x/n)^n$ gives

\[
 \sum_{p=3,5,\ldots} p|z_p|^{n}\le
 \sum_{p=3,5,\ldots}p(1+(p^2-1)/100)^{-50}
 \le 3(27/25)^{-50}+\tfrac{25}{49}(27/25)^{-49}<1/10.
\]

The summand decreases for $p\ge3$, so after its first term half the integral from $3$ to infinity bounds the rest. If $s\ge10$, $|T_s(w)|/|T_s(w_1)|<2$ for $|w|\le R$. The real part of the exponent in modulus is bounded above in absolute value by $e^{-s}(3s/2+2\pi+(3s/2+2\pi)^2/2)<1/20$, a decreasing bound from $10$. If $s\le10$, then $S^2>2/25$. Apply the preceding weighted-sum bound with exponent $n/2$ in place of $n$ (use $n/2\ge100$ and $(n/2)S^2\ge1$). The remaining exponent supplies a factor at most $(1+16/25)^{-n/4}$. Even after multiplication by $4e^{14}$ the weighted sum is less than $1/5$. Thus in both ranges the sum with the amplitude ratios is below $1/5$.

Crucially, the relative error is small enough for consecutive rows even where $S$ is small. Exactly

\[
 |z_p-\cos\theta|=S\frac{\sqrt{S^2+p^2\cos^2\theta}}{\sqrt{\cos^2\theta+p^2S^2}}\le pS\qquad(p>0).
\]

The difference of the principal contributions is $\cos(\Psi_n+\theta)-\cos\theta\cos\Psi_n=-S\sin\Psi_n$. The other poles have the additional factor just bounded. The circle estimates consequently prove

\[
 |Y_n-\cos\Psi_n|<1/4,\qquad
 |Y_{n+1}-\cos\theta\,Y_n+S\sin\Psi_n|<S/4. \tag{1}
\]

At $k=3$ these hold for every $n\ge400$ and $nS^2\ge1$.

For every fixed $k\ge3$ the same inequalities (1) hold on $nS^2\ge1$ for all sufficiently large $n$. To see this, the exponent of $T_s$ on $|w|\le R$ is $O_k(e^{-s}(1+s)^{k-1})$ as $s\to\infty$. On the complementary compact $s$-interval this exponent is bounded and $|T_s(w_1)|$ is bounded below. Thus the circle errors are $O_k(\sqrt n(17/20)^{n/2})=o(S)$. Above some fixed $s_k$ all amplitude ratios are below 2. Below $s_k$ the positive lower bound for $S$ supplies an exponentially small factor as in the $s\le 10$ argument.

\section{Coefficients at the overlap}

Cut the cycles at every ascent into $p$ labeled decreasing runs. A color word on $[n]$ assigns the labels to the runs, and a permutation of the runs specifies the upward boundary edges. Acceptance requires nonempty run classes, each boundary upward, and the total label count in each run-cycle at least $k$. Every admissible permutation has exactly $p!$ choices of run labels, canceled by the $p!$ possible uniform permutations of runs. Thus $\alpha_{n,p}=a_{n,p}/p^n$ is the acceptance probability, with $0\le\alpha_{n,p}\le1$. The run encoding is the one in the \href{https://github.com/ProgramComputer/open-math-research/pull/7}{reviewed extreme-rank proof}.

For $n\ge\max(400,25k^2)$ and $p\le(3/5)\sqrt n$, require each color to occur at least $k$ times in the full word and at least once in each half. Then every run permutation is accepted. The full-word count of one color is binomial of mean $\mu=n/p\ge(5/3)\sqrt n>8k$. Its probability below $k$ is at most $\exp(-49\mu/128)$ by the lower-tail bound. The shorter half has mean at least $(4/5)\sqrt n$, so its empty-color probability is at most $\exp(-4\sqrt n/5)$. The union bound gives

\[
1-\epsilon_n\le\alpha_{n,p}\le1,\qquad
\epsilon_n=\frac95\sqrt n e^{-3\sqrt n/5},\qquad
\sqrt n\,\epsilon_n<\frac1{100}.
\tag{2}
\]

The last expression decreases for $n\ge400$ and satisfies the bound at $n=400$ since $e>13/5$. The same interval contains $\alpha_{n+1,p}$ for these $p$, as $\epsilon_{n+1}\le\epsilon_n$.

\section{Exact first ranks and their consecutive signs}

For $1\le h\le(3/8)\sqrt n$, let

\[
u_{n,h}=\left(\frac h{h+1}\right)^n,\qquad
F_{n,h}(z)=\frac{E_{n,k}(u_{n,h}z)}{h^n u_{n,h}^h z^h}
=\sum_{p\ge1}\alpha_{n,p}
\left[\frac ph\left(\frac h{h+1}\right)^{p-h}\right]^n z^{p-h}.
\]

Coefficients beyond the degree are zero. Both bases at $p=h,h+1$ equal one. In the normalized bracket raised to $n$, successive terms away from this pair in either direction decrease by a ratio at most $q=e^{-5}<1/100$. Indeed the maximal ratios on the left and right are $(1-1/h^2)^n$ and $(1-1/(h+1)^2)^n$, respectively, and $n/(h+1)^2>5$ when $n\ge400$. The first ratio is unnecessary when $h=1$.

On $1/2\le|z|\le3$,

\[
|F_{n,h}(z)-(1+z)|\le4\epsilon_n+
\frac{2q}{1-2q}+\frac{9q}{1-3q}<\frac12 .
\]

On the circle $|z+1|=1/8$ the same bound with $|z|\le2$ gives $3\epsilon_n+6q/(1-2q)<1/8$. Rouché on the circles of radii $1/2$ and $3$ after multiplying by $z^{h-1}$ gives respectively $h-1$ and $h$ zeros of $P_n(u_{n,h}z)$. Rouché on the small offset circle gives a simple real negative root, since the disk is conjugation invariant. This establishes its absolute rank $h$, and its modulus belongs to $(7u_{n,h}/8,9u_{n,h}/8)$.

At such a root $t=u_{n,h}z$, consider $E_{n+1,k}(t)-hE_{n,k}(t)$ divided by $h^n u_{n,h}^h z^h$. The coefficient of each normalized term is then

\[
(p-h)\alpha_{n,p}+p(\alpha_{n+1,p}-\alpha_{n,p}).
\]

The $p-h$ part vanishes at $p=h$ and at $p=h+1$ contributes $\alpha_{n,h+1}z$, strictly below $-(7/8)(1-1/100)$. With $s=1/50$, the sum of absolute values of its remaining terms is at most

\[
\sum_{\ell\ge1}\ell s^\ell+2\sum_{\ell\ge1}(\ell+1)s^\ell
=\frac{3s}{(1-s)^2}+\frac{2s}{1-s}<\frac{11}{100}.
\]

For $p\le(3/5)\sqrt n$, (2) bounds the coefficient increment by $p\epsilon_n<3/500$. The absolute sum of the normalized bracket terms is less than four, so this portion contributes less than $3/125$. For $p>(3/5)\sqrt n$, one has $p-h>(9/40)\sqrt n\ge4.5$ and $p<3(p-h)$. The increment bound $p$ then gives a total less than $6\sum_{\ell\ge4}(\ell+1)s^\ell<1/1000$. The result is therefore strictly negative. Restoring the normalization, $E_{n+1,k}$ at the old root has sign $(-1)^{h+1}$ and $P_{n+1}$ has sign $(-1)^h$.

At the old root the modulus divided by $u_{n+1,h}$ lies in $(7/8,9/4)$, since $u_{n+1,h}/u_{n,h}=h/(h+1)$. This lies strictly inside the new annulus $(1/2,3)$, which has only its rank-$h$ root and exactly $h-1$ simple negative roots inside its lower boundary. Therefore the new sign puts the old rank $h$ strictly after the new rank $h$ and strictly before all larger new ranks.


\section{Phase checkpoints}

At $k=3$ the phase correction is $\gamma=\pi(a+s b)$. Its derivative in $\theta$ is $(s^2+\pi^2)s b(1-2a)$, which lies between zero and $27/e^3+\pi^2/e<6$. Its limits at $\theta=0,\pi/2$ are $0,\pi/2$. Thus $\Psi_n$ increases strictly from 0 to $n\pi/2$. For any fixed $k$, the correction still has bounded derivative. At $s\to\infty$ this follows by differentiating its finite polynomial expression and using $ds/d\theta=-(s^2+\pi^2)/\pi$. At $s=0$ all the odd powers of degree $l\ge3$ have coefficient zero because $A_{l-1}$ has odd degree and is reciprocal. The even powers give real terms and the linear term has imaginary part $\pi/2$. Hence the same phase endpoints and eventual strict increase hold at fixed $k$.

Write $H=\lfloor3\sqrt n/8\rfloor$. Let $\theta_j$ be the hit of $j\pi$ by $\Psi_n$ for $H+1\le j\le n/2$. Extend the samples to $j>n/2$ by $\theta_j=\pi-\theta_{n-j}$. In the lower half every such sample and every larger angle up to $\pi/2$ satisfies $nS^2\ge1$. At $k=3$, $\theta_{H+1}\ge(H+1)\pi/(n+1)>3\pi\sqrt n/(8(n+1))$; the sine of the last quantity exceeds $1/\sqrt n$ for $n\ge400$, by $\sin x\ge x-x^3/6$. At fixed $k$ the bounded phase correction gives the same eventually, since $3\pi/8>1$. By (1), the sign at $t_j=-\exp(-\pi\cot\theta_j)$ is $(-1)^j$, extended by reciprocity to the upper half of the angle range.

The preceding small-rank argument places the $H$ innermost roots and establishes the next-row signs there, uniformly for $n\ge\max(400,25k^2)$. Each rank $h\le H$ lies at modulus $x$ with $7/8<x/u_h<9/8$, where $u_h=(h/(h+1))^n$. In particular all $H$ precede the first checkpoint in increasing modulus. Indeed at $k=3$ its $s$ is at most $(n+1)/(H+1)$. At fixed $k$ it is at most $(n+1)/(H+1)+o(1)$ because $s\asymp\sqrt n$ there and $\gamma=O_k(e^{-s}(1+s)^{k-1})$. Finally

\[
 n\log(1+1/H)-(n+1)/(H+1)\ge n/(2(H+1)^2)-1/(H+1)>\log(9/8),
\]

using $H+1\le(17/40)\sqrt n$ for $n\ge400$ and the positive Taylor series of $-\log(1-1/(H+1))$.

Between the checkpoints $j=H+1,\ldots,n-H-1$ are $n-2H-2$ sign changes. The $H$ inner roots and their $H$ reciprocal partners supply the remaining $2H$. Degree exhausts all the intervals and gives simplicity. At a middle root with $h$th modulus on the lower half, $\Psi_n$ lies strictly between $h\pi$ and $(h+1)\pi$. For odd $n$ the root in the central interval is exactly $-1$ by reciprocity, where $\Psi_n=n\pi/2$ gives the same phase inequality. The pole comparison (1) gives $|\cos\Psi_n|<1/4$ and therefore $Y_{n+1}$ has sign $-\operatorname{sgn}\sin\Psi_n=(-1)^{h+1}$. On reciprocating a root the rank is $n-1-h$, and $E_{n+1,k}(1/t)=t^{-(n+1)}E_{n+1,k}(t)$ supplies the corresponding sign there. At every old nonzero root the sign of $E_{n+1,k}(t)/t$ is thus $(-1)^h$. Its value at zero is one and the new degree is one higher. The intermediate value theorem gives strict interlacing in all the asserted rows.

\section{Exact finite join at minimum length three}

Let $E_n=E_{n,3}$, with $E_0=1$ and $E_1=E_2=0$. Delete the largest label. If its cycle had length at least four, removal gives a permitted permutation on $n-1$ labels. Conversely insert the largest label in any of the $n-1$ directed edges of that permutation. Inserting into one of its $h$ ascending edges preserves the number of ascents and inserting into any other edge adds one. If its cycle had length three, choosing the two other labels and choosing one of the two cyclic orders gives weights $t$ and $t^2$. This proves

\[
 E_n=(n-1)tE_{n-1}+t(1-t)E'_{n-1}
       +\binom{n-1}{2}(t+t^2)E_{n-3}.
\]

An independent coefficient reconstruction begins with all permutations, uses inclusion-exclusion on fixed points to obtain $D_n=\sum_{j=0}^n(-1)^j\binom nj A_{n-j}$, and then uses inclusion-exclusion on transpositions:

\[
 E_n=\sum_{j=0}^{\lfloor n/2\rfloor}(-1)^j
       \frac{n!}{(n-2j)!2^j j!}\,t^j D_{n-2j}.
\]

Deleting the chosen fixed points or transpositions preserves the relative order on the other labels and their excedance statistic. A transposition contributes exactly one excedance. The second computation uses the ordinary Eulerian recurrence.

The archived records give the full vector for every $3\le n\le401$, together with two dyadic numbers in $(0,1)$ bracketing each root modulus below one. To evaluate $E_n(-x)/(-x)$ at such an endpoint the verifier uses integer Horner evaluation at fixed binary precision. Each rounded step introduces absolute error less than one scaled unit, and multiplication by $x<1$ contracts prior errors. A magnitude greater than the number of coefficients certifies the exact sign. The checker validates opposite signs, disjoint intervals within each row, and strict alternation of the intervals from consecutive rows. Reciprocity gives the partners above one and, in odd degree, the zero at $-1$. The number of distinct roots obtained equals the degree. These checks cover 399 rows, 398 consecutive pairs and 39,800 dyadic brackets. The analytic proof covers every pair beginning at $n\ge400$.


\section{A quadratic bound uniform in the minimum cycle length}


\textbf{Theorem 2.} An absolute $B$ suffices for the conclusions of the fixed-threshold theorem whenever $k\ge3$ and $n\ge Bk^2$. This includes simultaneous strict interlacing as $k$ grows up to an absolute multiple of $\sqrt n$. All comparisons are between sizes $n,n+1$ at the same $k$. We use the notation and normalization in the proof of Theorem 1.

Write $L(w)=\log T_s(w)$ with its polynomial branch. For $s>0$ the Eulerian power-sum identity gives

\[
 L(w)=aw-\sum_{m\ge1}\frac{(-1)^m e^{-ms}}{m}
                       \sum_{j=2}^{k-1}\frac{(mw)^j}{j!}.
 \tag{3}
\]

For every $s\ge0$ we also have

\[
 L(w)=aw-\sum_{j=2}^{k-1}\frac1j\sum_{p\ {\rm odd}}(w/w_p)^j.
 \tag{4}
\]

For (4), factor $1+e^{w-s}=2e^{(w-s)/2}\cosh((w-s)/2)$. The order-one Hadamard factorization centered at $w=0$ has genus at most one. Its logarithm has coefficient $-\sum_p w_p^{-j}/j$ in degrees $j\ge2$, by absolute convergence. Comparison with the Eulerian EGF gives (4). At $s=0$ the factorization is regular at zero.

For $s\ge10k$, keep the fixed-threshold proof's original circle, which has $R\le2s$. In (3) the truncated exponential with absolute argument $mR$ is at most $4^k\exp(mR/4)$. The whole nonlinear part is $O(4^k e^{-s/2})=O(e^{-s/3})$ throughout the circle; the linear part is $O(se^{-s})$. Along $w_1=s+i\pi$, differentiation of (3) costs at most a factor $m$ in each summand and a second copy of a truncated exponential. It gives the same exponential bound for the total derivative with respect to $s$. In particular the phase and its angular derivative are small there. For $s$ above an absolute constant the pole amplitude ratios are below two on this circle. If necessary increase $B$ to treat finitely many bounded $k,s$ using the compact-interval fixed-threshold argument.

For $0\le s\le10k$ choose a radius $R$ in $[r(1+1/k),r(1+2/k)]$ at radial distance at least $c/k^2$ from every pole. Indeed there are $O(1+s)=O(k)$ distinct pole radii below $3r$ and the interval has length $r/k$. Removing intervals of radius $c/k^2$ around them leaves an admissible choice for an absolute $c>0$. On the circle the reciprocal denominator is $O(k^2)$, by the elementary distance estimate in the fixed-threshold note (with a variable distance). We give three bounds for $L$.

1. For $0\le s\le\sqrt k$ and $|w|\le R$, equation (4) gives $|L(w)|\le C(1+s)\log(k+1)$. Indeed $|w/r|^j\le e^2$ and $\sum_{p\ {\rm odd}} |r/w_p|^j\le \sum_p |r/w_p|^2=O(1+s)$ for $j\ge2$, by comparison with the integral of $r^2/(s^2+\pi^2 p^2)$ away from the first poles. The term $aR$ is bounded. For the angular derivative at $w_1$ set $v_p=w_1/w_p=e^{i\theta}/(\cos\theta+ip\sin\theta)$. Then $|v'_p|=|p-1||v_p|^2$. Thus $\sum_p |v_p|^{j-1}|v'_p|\le\sum_p|p-1||v_p|^3=O((1+s)^2)$ by comparison with the integral or splitting at $|p|=1+s$. Summing over $j$ gives $|\gamma'(\theta)|=O(k(1+s)^2)=O(k^2)$. The linear term has a bounded angular phase derivative since $\Im(aw_1)=\pi a$.

2. For $\sqrt k\le s\le10k$ and $|w|\le R$ we have $|L(w)|\le C(1+\log^+(k/s))$. Here $R-s\le C_1(1/s+s/k)$ and $R/s$ is bounded by an absolute constant. For $m\le M=\lceil 10k/s\rceil$, $m(R-s)\le C_2$, so the absolute value of the $m$th summand of (3) is at most $e^{C_2}/m$. For the remaining terms and sufficiently large $k$, $R/s\le3/2$. Bound the truncated exponential by $2^k e^{mR/2}$. Then $e^{-ms}2^ke^{mR/2}\le e^{-ms/8}$ since $ms>10k$. The geometric tail is bounded uniformly. Finitely many smaller $k$ can again be handled by the fixed-threshold compact-interval argument. Differentiating along $w_1$ costs a factor at most $2m$ before taking the same absolute majorants (one sum runs only to $j=k-2$). Hence $|dL(w_1)/ds|=O(1+k/s)$, including the differentiated linear term. As $|ds/d\theta|=(s^2+\pi^2)/\pi$, this gives $|\gamma'(\theta)|=O(k^2)$.

3. The circle integral in this range divided by its principal amplitude is bounded by $C k^2 r\exp(2\sup |L|)(1+1/k)^{-n}$. Dividing by $S=\pi/r$ costs at most $O(k)$ more. Its logarithm is bounded by $C(1+\sqrt k)\log(k+1)+O(\log k)-n/(2k)$. Increasing $B$ gives $|e_n|+|e_{n+1}|<S/40$ for $n\ge Bk^2$. For other included poles, split $|z_p|^{n+1}$ at exponent $n/2$. Since $s\le10k$ and $n\ge Bk^2$, we have $(n/2)S^2\ge c_0 B\ge1$ after increasing $B$. The fixed-threshold weighted integral bound applies to the first factor. The second is at most $(1+8S^2)^{-n/4}\le\exp(-c_1 n/(1+s^2))$. In $s\le\sqrt{k}$ this exponential absorbs the bound $\exp(C(1+\sqrt k)\log(k+1))$ for the amplitude ratio, uniformly by enlarging $B$. In $s\ge\sqrt k$ it absorbs $\exp(C(1+\log^+(k/s)))$ since $\sup_{x\ge1/10}(1+\log^+ x)/x^2<\infty$. This proves both principal-pole estimates whenever $nS^2\ge1$. In the range $s\ge10k$ these estimates already follow from the original fixed-threshold circle and the first bound after (4).

By the phase derivative bound $\Psi_n$ increases for $n\ge Bk^2$ with the same endpoints. At $s=10k$ its phase is at least $(n+1)\arctan(\pi/(10k))-O(1)$, greater than $(H+1)\pi$ for $H=\lfloor3\sqrt n/8\rfloor$ by choosing $B$ large. Thus the first checkpoint occurs where $s>10k$. At that checkpoint $\theta=\Theta(n^{-1/2})$ because $\gamma$ is bounded there, and consequently $s=\Theta(\sqrt n)$. Its phase correction is then $O(e^{-c\sqrt n})$ uniformly. The lower sine bound, the overlap with the $H$ innermost roots, the exact rank count and the consecutive sign comparison follow from the fixed-threshold proof, with the constant gap between $3\pi/8$ and one. The extreme-rank lemma already holds for every $n\ge\max(400,25k^2)$. This proves uniformity in $k$. The proof establishes existence of $B$; the explicit cutoff 400 at $k=3$ is obtained separately.
\section{Consecutive interlacing through the two-cycle range}
\label{sec:two-cycle-insertion}

\textbf{Two-cycle interlacing theorem.}
Let $k\ge3$. For every $k\le n\le3k-2$, put $P_n(t)=E_{n,k}(t)/t$.
Then $P_n$ and $P_{n+1}$ have simple negative zeros, and the zeros of
$P_{n+1}$ strictly interlace those of $P_n$ with the new polynomial
supplying both outer zeros. More precisely, at every zero $x$ of $P_n$,
\begin{equation}
 \left|L(x)\right|<\left|D(x)\right|,\qquad
 D(t)=t(1-t)P_n'(t),\qquad
 L(t)=\binom{n}{k-1}A_{k-1}(t)E_{n+1-k,k}(t).
 \label{eq:two-cycle-lag}
\end{equation}
Thus the conclusion covers every consecutive pair for which the new row
admits at most two cycles.

To obtain the insertion recurrence, either insert the largest label
into an edge of an existing cycle or place it in a new $k$-cycle.
Inserting it after the initial label of an ascending edge preserves
the excedance count. Inserting it in any other edge adds one
excedance. With the convention $E_{0,k}=1$ and $E_{m,k}=0$ for
$1\le m<k$, this gives
\begin{equation}
 P_{n+1}(t)=\bigl(1+(n-1)t\bigr)P_n(t)+D(t)+L(t).
 \label{eq:two-cycle-recurrence}
\end{equation}
The two polynomials $P_{n+1}$ and $P_{n+1}-2L$ therefore take the values
$D+L$ and $D-L$, respectively, at an old root. For $k\le n\le2k-2$
only single cycles occur in both rows and $L=0$. Eulerian
interlacing follows from
$A_{m+1}=(1+mt)A_m+t(1-t)A_m'$ by induction from $A_1=1$,
evaluating the recurrence at the simple negative zeros of $A_m$.
The first nonzero insertion term, at $n=2k-1$, is covered by the
\href{https://github.com/ProgramComputer/open-math-research/pull/27}{reviewed first-transition argument}.
It compares $A_{2k-2}$ with both polynomials
$A_{2k-1}\pm\binom{2k-1}{k-1}tA_{k-1}^2$.
We treat $2k\le n\le3k-2$ below.

\subsection{The first ranks and their derivative bound}

Write $M=n-1$, $E=E_{n,k}$, and
$\alpha_p=[t^p]E/p^M$ for $p\ge1$. We take $n\ge100$ in this
subsection. The bounds needed for the first ranks are
\begin{equation}
 0\le\alpha_p\le3\quad(p\ge1),\qquad
 \alpha_p\ge99/100\quad
 \left(p\le \frac{451}{1000}\sqrt M\right).
 \label{eq:two-cycle-coefficients}
\end{equation}
A permutation with $p-1$ descents gives a word on $p$ colors by assigning
each label to its increasing run. The word determines the permutation,
so $[t^p](tA_m)\le p^m$. Conversely, sorting each color class and
concatenating the classes in color order creates all $p-1$ descents
if every class meets both label halves. The union bound shows that
the single-cycle term $tA_M$ alone gives
$\alpha_p\ge 1-2p\exp(-\lfloor M/2\rfloor/p)$. For the stated
range of $p$, this error is below
$\sqrt M\exp(-9\sqrt M/10)<1/100$.

Each unordered pair of cycles is counted twice in the sum
\begin{equation}
 E(t)=tA_{n-1}(t)+\frac12\sum_{j=k}^{n-k}\binom nj
                (tA_{j-1}(t))(tA_{n-j-1}(t)).
 \label{eq:two-cycle-exact-polynomial}
\end{equation}
Consequently the coefficient bound from words, followed by symmetry
under $a\mapsto p-a$, gives
\[
 \alpha_p\le
 1+\frac12\sum_{a=1}^{p-1}\frac{p}{a(p-a)}
        \Pr\{\operatorname{Bin}(n,a/p)\in[k,n-k]\}
 =
 1+\sum_{a=1}^{p-1}\frac1a
        \Pr\{\operatorname{Bin}(n,a/p)\in[k,n-k]\}.
\]
Since $k\ge n/3$, Hoeffding's inequality bounds the probability by
$\exp(-8n/225)$ whenever $a/p\notin(1/5,4/5)$. For $p\le n-1$ the sum of
these terms is at most $(1+\log n)\exp(-8n/225)<1/6$.
The expression decreases for $n\ge100$ and satisfies the bound at
$n=100$ using $\log100<14/3$ and $\exp(32/9)>34$.
For the remaining indices the sum of $1/a$ is at most $11/6$.
When $p\le4$, this follows by summing through $a=3$.
When $5\le p\le11$, all the indices lie in $[2,8]$.
For $p\ge12$, integrate $1/a$ from the first index to the last,
allowing the first term separately, to obtain
$\log4+5/p\le\log4+5/12<11/6$. Coefficients beyond $n-1$ vanish.
This proves \eqref{eq:two-cycle-coefficients}.

Put $H=\lfloor(7/20)\sqrt M\rfloor$.
For a rank $1\le h\le H$ define
\[
 \rho=\frac{\alpha_h}{\alpha_{h+1}},\qquad
 u=\rho\left(\frac h{h+1}\right)^M,\qquad
 G(z)=\frac{E(uz)}{\alpha_h h^M(uz)^h},\qquad
 b_p=\left[\frac ph\left(\frac h{h+1}\right)^{p-h}\right]^M.
\]
We have $h+1\le(451/1000)\sqrt M$ and
$\rho,\rho^{-1}\le R=100/33$. The terms of $G$ with indices
$h,h+1$ are exactly $1,z$. The successive ratios of $b_p$ away from
this pair are at most $(1-1/h^2)^M$ on the left and
$(1-1/(h+1)^2)^M$ on the right. Both are bounded by
$q_0=\exp(-1000000/203401)<1/135$, and the left side is empty
when $h=1$.

Whenever $q_0\max(|\rho z|,|\rho z|^{-1})\le\eta$, the absolute values
of the terms of $G$ at indices $h-\ell$ and $h+1+\ell$ are at most
$R\eta^\ell$ and $R|z|\eta^\ell$, respectively, for $\ell\ge1$.
The second estimate uses the coefficient ratio to $\alpha_{h+1}$,
as the right central term equals $z$.
On $|z+1|=1/4$ we may take $\eta=3/100$, since
$3/4\le |z|\le5/4$. Here
\[
 |G(z)-(1+z)|\le R(1+|z|)\frac{\eta}{1-\eta}<1/4.
\]
On $|z|=1/2$ and $|z|=3$ we may take $\eta=7/100$, and the same
bound is below $1/2$ and $2$, respectively. Rouch\'e's theorem
applied to $z^{h-1}G(z)$ gives $h-1$ roots inside $|z|=1/2$, $h$
inside $|z|=3$, and one in $|z+1|<1/4$. That disk lies in the
annulus between the two circles. Its sole root is real by
conjugation. Thus $P_n$ has a simple negative root $x=uz$ of absolute
rank $h$ with $3/4<|x|/u<5/4$.

At this root, put $X=|x|$. The normalized derivative
$xE'(x)/(\alpha_h h^M x^h)$ equals $zG'(z)-G(z)$ because
$G(z)=0$. Its central terms contribute $-1$ together. The remaining
weights are $\ell+1$ and $\ell$ on the left and right. With
$\eta=3/100$ this gives
\begin{equation}
 \begin{split}
 \bigl|zG'(z)-G(z)+1\bigr|
 &\le R\left((1+5/4)\frac{\eta}{(1-\eta)^2}
                    +\frac{\eta}{1-\eta}\right)<\frac13,\\
 |xE'(x)|&>\frac23\alpha_h h^M X^h .
 \end{split}
 \label{eq:two-cycle-front-derivative}
\end{equation}

The residual row has size $n+1-k$ between $k$ and $2k-1$.
It consists of single cycles, so
$L(t)=\binom{n}{k-1}tA_{k-1}(t)A_{n-k}(t)$. Applying the same
coefficient bound to both Eulerian factors yields
\[
 [t^r](tL)
 \le \frac{n}{n-k+1}\,r^M
     \sum_{a=1}^{r-1}
         \Pr\{\operatorname{Bin}(M,a/r)=k-1\}.
\]
Here $b=(k-1)/M$ belongs to $[1/3,1/2]$, and Stirling's bounds give
\[
 \Pr\{\operatorname{Bin}(M,b)=k-1\}
 \le \frac{\exp(1/(12M))}{\sqrt{2\pi M b(1-b)}}
 < \frac{7}{8\sqrt M}.
\]
Changing the binomial parameter from $b$ to $y$ multiplies this
point probability by $\exp(-M I(b,y))$, where
$I(b,y)=b\log(b/y)+(1-b)\log((1-b)/(1-y))$.
As a function of $b$, $I$ and its first derivative vanish at $b=y$,
and its second derivative is at least four. Thus
$I(b,y)\ge2(b-y)^2$. The Gaussian sum on the lattice $y=a/r$ is
at most $1+\sqrt{\pi/2}\,r/\sqrt M$. Indeed, each superlevel
interval contains at most its length plus one lattice points in
the $a$ coordinate, and integrating the level from zero to one
gives the bound. Since $n/(n-k+1)<2$,
\begin{equation}
 [t^r](tL)\le \frac{2(7/8)}{\sqrt M}\,
       r^M\left(1+\sqrt{\pi/2}\,\frac r{\sqrt M}\right).
 \label{eq:two-cycle-front-lag}
\end{equation}
For the sum at $x$, the weights
$r^M X^r/(\alpha_h h^M X^h)$ are bounded by
$(100/99)\eta^\ell$ at $r=h-\ell$ and
$(100/99)(5/4)\eta^\ell$ at $r=h+1+\ell$, for $\ell\ge0$.
For the right-hand weights use $1/\alpha_{h+1}$, which absorbs
the factor $\rho$ from the first step. The weight in parentheses in
\eqref{eq:two-cycle-front-lag} is at most $157/100$ through
$r=h+1$ and grows by less than $13/100$ per step thereafter. The
normalized weighted sum is therefore below
\[
 \frac{100}{99}\left(\frac94\frac{157}{100(1-\eta)}
          +\frac54\frac{13\eta}{100(1-\eta)^2}\right)
 <\frac{37}{10}.
\]
Since $E'(x)=xP_n'(x)$, equations
\eqref{eq:two-cycle-front-derivative} and
\eqref{eq:two-cycle-front-lag} give
\begin{equation}
 \frac{|L(x)|}{|D(x)|}
 =\frac{|xL(x)|}{(1+X)|xE'(x)|}
 <\frac{777}{80\sqrt M}<1.
 \label{eq:two-cycle-front-ratio}
\end{equation}

\subsection{Middle-root signs and a fixed rational comparison}

For $t=-e^{-s}$, $s\ge0$, let
\[
 r=\sqrt{s^2+\pi^2},\qquad
 \theta=\arg(s+i\pi),\qquad S=\sin\theta,\qquad
 z_p=(\cos\theta+ipS)^{-1},\qquad
 B_m=\sum_{p\ \mathrm{odd}}z_p^m .
\]
Here $p$ runs over both positive and negative odd integers.
The Eulerian power-sum identity gives the normalized $m$-cycle weight
\begin{equation}
 \frac{r^m tA_{m-1}(t)}{m!(1-t)^m}=\frac{B_m}{m}\qquad(m\ge2).
 \label{eq:two-cycle-pole}
\end{equation}
Indeed $tA_{m-1}(t)/((m-1)!(1-t)^m)$ is the coefficient of $w^{m-1}$
in $1/(1+e^{w-s})-1$. The identity at $t=-1$ follows by continuity.
The poles $s+ip\pi$ all have residue $-1$.
For fixed $s$, the circles of radii $\sqrt{s^2+4b^2\pi^2}$, with integers
$b\to\infty$, stay a positive distance from the poles. On these
circles the reciprocal denominator is bounded: if the real part
of $w-s$ is bounded away from zero the reverse triangle inequality
applies, and otherwise its imaginary part stays away from the odd
multiples of $\pi$. The coefficient integral on the circle tends to
zero. Summing the residues in the coefficient formula proves
\eqref{eq:two-cycle-pole}, since the sum converges absolutely.

Define
\[
 F_n=\frac{r^n E_{n,k}(t)}{n!(1-t)^n},\qquad
 f=\frac n2F_n
   =\frac{B_n}{2}+\sum_{j=k}^{n-k}w_j\frac{B_jB_{n-j}}{2},
 \qquad w_j=\frac{n}{2j(n-j)}.
\]
The identity follows by applying \eqref{eq:two-cycle-pole} to
\eqref{eq:two-cycle-exact-polynomial}. Set $W=\sum_{j=k}^{n-k}w_j$.
Then $W=\sum_{j=k}^{n-k}1/j\le\log 2<7/10$, by integration from
$k-1$ and the bound $n-k\le2k-2$.

Assume $n\ge N=192$ and $nS^2\ge c=6/5$.
Let $\beta=c/N=1/160$ and $K(p)=1+\beta(p^2-1)$. If $m/n\ge a$,
the monotonicity of $(1+u/n)^n$ in $n$ for $u\ge0$ gives
$|z_p|^m\le K(p)^{-H_a}$, where $H_a=Na/2$.
All the parameters $a$ used below are at least $1/3$ and are integer
multiples of $1/48$, so $H_a$ is an integer. With
\[
 I_a=\frac{K(9)^{1-H_a}}{4\beta(H_a-1)},\qquad
 T(a)=\sum_{p=3,5,7,9} K(p)^{-H_a}+\frac{I_a}{9},\qquad
 U(a)=\sum_{p=3,5,7,9}pK(p)^{-H_a}+I_a ,
\]
we obtain, for $e_m=B_m/2-\cos(m\theta)$,
\begin{equation}
 |e_m|\le T(a),\qquad |e_m'|\le mU(a).
 \label{eq:two-cycle-single-tails}
\end{equation}
The term with $|p|=1$ gives the displayed cosine. For the derivative,
$|z_p'|\le |p|\,|z_p|^2$ and $|z_p|\le1$. The sum on positive odds
above nine is at most half the integral starting at nine of the
decreasing majorant. The weighted integrand $pK(p)^{-H_a}$ decreases
there because $(2H_a-1)\beta p^2>1-\beta$. Its half integral equals
$I_a$, and the unweighted integrand is at most one ninth of the weighted
one on that interval. These summable bounds also justify
differentiating the series in this region.

Suppose $u\le\min(j,n-j)/n\le v$ for one adjacent mesh interval
$[u,v]$. Put
\begin{align*}
 A_{uv}&=1+2\bigl(T(u)+T(1-v)+T(u)T(1-v)\bigr),\\
 C_{uv}&=1-2u+
 2v\bigl(T(1-v)+U(u)(1+T(1-v))\bigr)\\
 &\hspace{2.65em}
 +2(1-u)\bigl(T(u)+U(1-v)(1+T(u))\bigr).
\end{align*}
For the two principal cosines, the identity
$2\cos(j\theta)\cos((n-j)\theta)=\cos(n\theta)+\cos((2j-n)\theta)$
gives a residual bounded by one in magnitude and $n(1-2u)$ in
derivative magnitude. Applying \eqref{eq:two-cycle-single-tails} and
the product rule to the terms involving at least one $e_m$ bounds
the full residual by $A_{uv}$ and its derivative by $nC_{uv}$.

For the threshold interval with index $i\in\{16,\ldots,23\}$ choose
$q=i/48\le k/n\le q_+=(i+1)/48$. Let $A_i,C_i$ be the respective
maxima of $A_{uv},C_{uv}$ over the adjacent mesh intervals in
$[q,1/2]$. With $T_*=T(1)$ and $U_*=U(1)$ this proves
\begin{equation}
 \begin{split}
 |f-(1+W)\cos(n\theta)|&\le T_*+WA_i,\\
 |f'/n+(1+W)\sin(n\theta)|&\le U_*+WC_i.
 \end{split}
 \label{eq:two-cycle-mixed-tails}
\end{equation}

To compare the insertion terms at a root $t=x$, multiply $D(x)$
and $L(x)$ by $x r^{n+1}/(n!(1-x)^{n+1})$.
Their transformed values are, respectively,
$S F_n'(\theta)$ and $B_kB_{n+1-k}/(n+1-k)$.
For the first equality use $t\,d\theta/dt=\pi/r^2$.
Terms from differentiating the normalization vanish at a zero of the
old row. For the second, apply
\eqref{eq:two-cycle-pole} to each of the two factors of $tL$.
Dividing further by $2S$ gives derivative magnitude $|f'|/n$
and lag magnitude at most
\[
 J_i=\frac{2(1+T(q))(1+T(1-q_+))}{15(1-q_+)},
\]
since $(n+1-k)S\ge(1-q_+)\sqrt{nc}>15(1-q_+)$.
The weight $W$ satisfies $0\le W\le W_i$ with
\[
 W_i=\min\left(\frac7{10},\frac{1-2q+1/N}{q}\right).
\]
The second bound uses at most $n-2k+1$ terms each at most $1/k$.

Define
\[
 R_i(w)=\frac{(T_*+wA_i)^2+(U_*+wC_i+J_i)^2}{(1+w)^2}.
\]
The table gives upper integer bounds for $10^6R_i(w)$ at the two
endpoints, computed by rational arithmetic from the formulas above.
The last two columns are ceilings.
\[
\begin{array}{c|c|r|r}
 i & W_i & \lceil 10^6R_i(0)\rceil & \lceil 10^6R_i(W_i)\rceil\\ \hline
 16 & 7/10  & 85473  & 988936\\
 17 & 7/10  & 88476  & 921937\\
 18 & 49/72 & 92155  & 844623\\
 19 & 41/76 & 96575  & 649634\\
 20 & 33/80 & 101824 & 488795\\
 21 & 25/84 & 108020 & 358384\\
 22 & 17/88 & 115311 & 256385\\
 23 & 9/92  & 123892 & 182547
\end{array}
\]
Since $A_i\ge1$, the polynomial
\[
 (T_*+wA_i)^2+(U_*+wC_i+J_i)^2-(1+w)^2
\]
is convex in $w$. Its endpoint values are strictly negative, so
\begin{equation}
 (T_*+WA_i)^2+(U_*+WC_i+J_i)^2<(1+W)^2.
 \label{eq:two-cycle-quadratic}
\end{equation}
At a checkpoint $\theta=l\pi/n$ in this region,
\eqref{eq:two-cycle-mixed-tails} and \eqref{eq:two-cycle-quadratic}
give $\operatorname{sgn} f=(-1)^l$.
At a zero of $f$, the first bound in \eqref{eq:two-cycle-mixed-tails}
bounds the cosine in the same inequality. Hence
$(1+W)|\sin(n\theta)|>U_*+WC_i+J_i$. The second bound in
\eqref{eq:two-cycle-mixed-tails} gives $|f'|/n>J_i$.
This proves \eqref{eq:two-cycle-lag} at all old roots in the
middle region on the side $-1\le t<0$.

\subsection{Joining the regions and counting signs}

Take $H=\lfloor(7/20)\sqrt{n-1}\rfloor$ as above and suppose $n\ge192$.
At the first checkpoint $\theta_{H+1}=(H+1)\pi/n$ we have
$n\sin^2\theta_{H+1}>6/5$. Indeed
$b=(7\pi/20)\sqrt{n-1}/n<\theta_{H+1}$ and
$b^2<121/(100n)$. Using $\sin b\ge b-b^3/6$ and
$333/106<\pi<22/7$ gives
\[
 n\sin^2\theta_{H+1}>
 \left(\frac7{20}\frac{333}{106}\right)^2
 \frac{191}{192}\left(1-\frac{121}{115200}\right)^2>\frac65.
\]
All the first $H$ roots precede this checkpoint in increasing modulus.
Their moduli are bounded by $(5/4)R(H/(H+1))^{n-1}$, and the
checkpoint corresponds to $s=\pi\cot\theta_{H+1}<n/(H+1)$.
With $M=n-1$, the positive Taylor series of
$-\log(1-1/(H+1))$ gives
\[
 M\log(1+1/H)-\frac{M+1}{H+1}
 \ge\frac{M}{2(H+1)^2}-\frac1{H+1}
 \ge\frac{1}{2(451/1000)^2}-\frac14
 >\log((5/4)R).
\]
This proves the strict separation.

Let $t_l=-\exp(-\pi\cot(l\pi/n))$ for
$l=H+1,\ldots,n-H-1$. The lower-half checkpoints through
$l=\lfloor n/2\rfloor$ satisfy the middle condition,
so $E_{n,k}(t_l)$ has sign $(-1)^l$ there. The upper-half
checkpoints follow by reciprocity. In fact
$t_{n-l}=1/t_l$ and $E_{n,k}(1/t)=t^{-n}E_{n,k}(t)$,
so the signs are again $(-1)^l$.
The $n-2H-2$ gaps between consecutive checkpoints each contain a
negative old root. The first $H$ roots and their $H$ reciprocal
partners lie strictly outside these gaps. This exhausts the degree
$n-2$ of $P_n$, so all old roots are simple.
For odd $n$ the center belongs to the central gap, and for even $n$
it is a checkpoint. Every old root on the lower side is therefore
covered by \eqref{eq:two-cycle-front-ratio} or
\eqref{eq:two-cycle-quadratic}. At an old root $x$, the reciprocal
identity for $P_n$ implies $D(1/x)=x^{1-n}D(x)$.
The cycle factors of $L$ give $L(1/x)=x^{1-n}L(x)$.
This transfers \eqref{eq:two-cycle-lag} to the upper side as well.

Since $P_n(0)=1$, the sign of $D$ at the old root of absolute rank
$h$ is $(-1)^h$. The two next-row polynomials $P_{n+1}$ and
$P_{n+1}-2L$ have constant and leading coefficients one, with
degree $n-1$. At every old root their signs equal the sign of $D$
by \eqref{eq:two-cycle-lag}. Both therefore have a root in each
successive gap and in each of the two outer intervals on the
negative axis. This exhausts their degree and proves strict
interlacing for $n\ge192$.

\subsection{Exact replay for the remaining pairs}

The remaining range has $3\le k\le95$ and
$2k\le n\le\min(191,3k-2)$. It contains $3\,007$ consecutive pairs
with $379\,502$ old roots, counting reciprocal partners and any
central roots. The finite verifier constructs the original coefficients
by distinguishing the cycle containing label one:
\[
 P_n(t)=A_{n-1}(t)+t\sum_{j=k}^{n-k}
              \binom{n-1}{j-1}A_{j-1}(t)A_{n-j-1}(t).
\]
It checks \eqref{eq:two-cycle-recurrence} separately as an integer
polynomial identity and checks both next-row polynomials against the
same old polynomial. Eulerian coefficients are reconstructed by their
recurrence and cross-checked by finite differences. A separate
Eulerian sign computation and the transition at $n=2k-1$ initialize the
brackets for each threshold.

Below modulus one, the verifier locates disjoint dyadic brackets
with denominator $2^{256}$. The signs of the old polynomial at the
bracket endpoints of rank $h$ are $(-1)^{h-1}$ and $(-1)^h$ in
order of increasing modulus. The brackets, their reciprocal partners,
and the center when the old degree is odd exhaust the old degree.
Each next-row polynomial is tested for sign $(-1)^h$ at both endpoints of that bracket. The sign
calculation uses scaled integer Horner evaluation. At each step the
rounding error is less than one scaled unit, and multiplication
by the modulus contracts the accumulated error. Hence an accumulator
with magnitude exceeding the degree in those units certifies the sign.
The endpoint signs at zero and $-1$ are evaluated exactly.

For each next-row polynomial, these signs force a lower root in
every gap between old brackets and one before the first bracket.
If the old degree is odd, the sign checked at its central root
forces an additional new root between the last lower bracket and
$-1$. If the new degree is odd, reciprocity gives its central root.
Reciprocal partners and these roots exhaust the new degree.
Thus each next-row polynomial has the same fixed sign throughout
each old bracket, and the central check and reciprocity give the
corresponding signs at the other old roots. This proves
\eqref{eq:two-cycle-lag} and strict interlacing for all finite pairs.

The \href{https://github.com/ProgramComputer/open-math-research/tree/research/topics/long-cycle-excedances}{topic archive}
contains the rational bound checker, the finite verifier and its replay
materials. Run the following from that topic directory with Python~3
and \texttt{python-flint} installed. FLINT is used for integer
polynomial arithmetic, and root signs use integer Horner evaluation.
\begin{verbatim}
python3 -I -O verify_two_cycle_bounds.py
python3 -I -O verify_two_cycle_band.py
\end{verbatim}
The second command reconstructs the brackets and needs no root archive.
Its optional \texttt{--records} argument writes a gzip JSONL index with
the certified pair ranges and hashes of the coefficients and endpoints.
The first command checks all sixteen rational endpoint comparisons and
the checkpoint join. No Lean or Rocq verification is asserted here.

\section{Effective front roots for all minimum cycle lengths}

\textbf{Partial theorem.} Let $n\ge10{,}000$ and $3\le k\le n$. Set $H=\lfloor7\sqrt n/20\rfloor$ and $P_n=E_{n,k}/t$. For every $1\le h\le H$, the root of $P_n$ of absolute rank $h$ is simple and negative. At each such root $x$, the distinguished-cycle term $L=\binom n{k-1}A_{k-1}E_{n+1-k,k}$ in the insertion recurrence satisfies
\[
|L(x)|<|x(1-x)P_n'(x)|.
\]
In order of increasing absolute value, the new rank-$h$ root is strictly closer to zero than the old rank-$h$ root. All higher-rank new roots are farther from zero than that old root. All $H$ old roots precede the phase sample $t=-\exp(-\pi\cot((H+1)\pi/n))$. The reciprocal conclusions hold at the other end. There is no restriction on the number of cycles. The proof compares the run-color probabilities for consecutive rows and bounds the remaining terms on three circles.

Use the \href{https://github.com/ProgramComputer/open-math-research/blob/main/topics/long-cycle-excedances/proof.md}{run encoding and insertion recurrence}. For each $p\ge1$ write $\alpha_{N,p}=[t^p]E_{N,k}/p^N$, which is the probability that an independent uniform word of length $N$ on $p$ run colors and a uniform permutation of the colors satisfy all upward boundaries and minimum lengths. Write $f(N,p,k)$ for the acceptance probability for the minimum lengths alone. It bounds $\alpha_{N,p}$ above. Put $f(N,p,k)=0$ for $0<N<k$ when $p\ge1$. The good event that every color appears in both label halves guarantees the upward boundaries and all nonempty colors, so
\[
0\le f(N,p,k)-\alpha_{N,p}\le 2p e^{-\lfloor N/2\rfloor/p}.
\tag{F1}
\]
This argument uses the probability form of the coefficients from the \href{https://github.com/ProgramComputer/open-math-research/blob/main/topics/long-cycle-excedances/uniform-interlacing.md}{uniform proof} and the adaptive root scaling from the \href{https://github.com/ProgramComputer/open-math-research/blob/main/topics/long-cycle-excedances/two-cycle-front.md}{two-cycle front estimate}.

\subsection{Effective run-color acceptance envelopes}

Put $J=\lfloor\sqrt n/2\rfloor$. For $p\le J$ and $a=kp/n$, the old row satisfies
\[
\frac{99}{100(1+(8/3)a)}\le\alpha_{n,p}\le
\min\left(1,\frac{151}{100a}\right).
\tag{F2}
\]
The upper bound also applies to $f(N,p,k)$ with $a$ replaced by $kp/N$ for every $n/2\le N\le n$ and $1\le p\le J$.

For the lower bound when $k\le n/2$, let $b=\lceil (4/3)a\rceil$, so $b\le p$. Write $d(p,b)$ for the probability that a uniform permutation of $p$ colors has no cycle shorter than $b$. Distinguishing the cycle containing a fixed color gives $1/(2b-1)\le d(p,b)\le1/b$ for $p\ge b$. Condition the run permutation to have no cycle shorter than $b$. A cycle of $u\ge b$ colors receives a binomial number of labels of mean $nu/p\ge4k/3$, so its probability of receiving fewer than $k$ labels is at most $e^{-nu/(32p)}$. For $u<p$, the unconditional expected count of marked $u$-cycles whose remaining cycles also have length at least $b$ is $d(p-u,b)/u\le1/(bu)$. A cycle of length $p$ receives every label and cannot fail. After division by $d(p,b)$, the union bound over the marked cycles is therefore less than
\[
2\sum_{u\ge1} \frac{e^{-nu/(32p)}}{u}
\le \frac{2e^{-25/4}}{1-e^{-25/4}}<\frac1{250}.
\]
The error in (F1) relative to $d(p,b)$ is bounded by $4p^2 e^{-(n-1)/(2p)}\le n e^{-(n-1)/\sqrt n}<1/1000$. Thus $\alpha_{n,p}\ge(99/100)d(p,b)\ge99/[100(1+(8/3)a)]$. When $k>n/2$, acceptance for the size condition requires a single run cycle and has probability $1/p$. Applying (F1) gives the same lower bound.

For the upper bound take $\mu=N/p\ge\sqrt n$, $a'=k/\mu$ and $b=\max(1,\lceil (2/3)a'\rceil)$. The bound is immediate if $k>N$ or $a'\le3/2$. Otherwise, acceptance with a run cycle of $u<b$ colors requires that this cycle receive at least $k$ labels despite mean $Nu/p<(2/3)k$. The exponential binomial tail is at most
\[
\exp[-k(\log(3/2)-1/3)]\le e^{-k/15}.
\]
Using the unconditional expected number $1/u$ of $u$-cycles gives
\[
f(N,p,k)\le d(p,b)+(1+\log p)e^{-k/15}
\le\frac{3}{2a'}+(1+\log p)e^{-a'\mu/15}.
\]
The last error times $a'$ is at most $(3/2)(1+\log(\sqrt n/2))e^{-\sqrt n/10}<1/100$. This uses decrease in $a'$ for $a'\ge3/2$, decrease in $\sqrt n$ for $\sqrt n\ge100$, and the bounds $p\le\sqrt n/2$ and $\mu\ge\sqrt n$. This proves (F2).

\subsection{Lag coefficients and consecutive row stability}

Write $m=n+1-k$. If $m<k$ the lag vanishes and we set $\lambda_p=0$. Otherwise $n/2\le m\le n$, and for $p\ge1$ we put
\[
\lambda_p=\sum_{j=1}^{p-1} b(j/p)f(m,p-j,k),
\qquad
b(y)=\binom n{k-1}y^{k-1}(1-y)^m.
\]
Then
\[
[t^p](tL)\le p^n\lambda_p,\qquad
0\le f(n+1,p,k)-f(n,p,k)\le \lambda_p/p .
\tag{F3}
\]
For the first bound, the coefficient of $t^j$ in $tA_{k-1}$ is at most $j^{k-1}$ by assigning the $k-1$ labels of a permutation to its increasing runs. Convolve with $[t^{p-j}]E_{m,k}=(p-j)^m\alpha_{m,p-j}$. Division by $p^n$ uses $k-1+m=n$, giving the one binomial point mass in (F3). For the second bound, couple the word of length $n+1$ to the old word by appending an independent color. If the minimum-size condition changes from false to true, the new label hits a run cycle that had precisely $k-1$ labels. A cycle of $j$ colors contributes expected count $1/j$ and is hit with probability $j/p$. Conditional on the count $k-1$, the remaining labels use the other $p-j$ colors independently and uniformly, and the remaining run permutation is uniform. The all-color cycle cannot have had size $k-1$. The same argument makes both differences zero when $m<k$, and both vanish for $p=1$. For use outside the range of (F2), $\lambda_p\le p$ follows directly.

For $p\le J$ the useful uniform estimate is
\[
\lambda_p\le \frac8{25}\alpha_{n,p}.
\tag{F4}
\]
If $k\le3n/(4p)$, set $\mu=n/p$. The sample $b(j/p)$ is bounded by $\Pr\{\mathrm{Bin}(n,j/p)\le k-1\}\le e^{-j\mu/32}$ for every $j\ge1$. Hence $\lambda_p<1/500$, while $\alpha_{n,p}\ge33/100$ by (F2).

Otherwise $k-1\ge150$. Put $y_0=(k-1)/n$. A sample with $j/p\le y_0+1/10$ has residual parameter
\[
\frac{k(p-j)}m\ge\frac45\frac{kp}n
\]
because $m/n=1-y_0\ge1/2$. The upper bound from (F2) for the residual and the lower bound for $\alpha_{n,p}$ give
\[
\frac{f(m,p-j,k)}{\alpha_{n,p}}
\le \frac{100}{99}(1+(8/3)a)\min\left(1,\frac{151}{80a}\right)
<\frac{61}{10}.
\]
For $j/p>y_0+1/10$, the ratio of $b(j/p)$ to its maximum at $y_0$ is $\exp(-nD(y_0\|j/p))$. Here $D$ is Bernoulli relative entropy and $D(v\|w)\ge2(v-w)^2$, by the second derivative in $v$. Hence $b(j/p)\le e^{-n/50}$. Summing these samples after division by $\alpha_{n,p}$ costs at most $n e^{-n/50}<1/10{,}000$, using $p(1+(8/3)a)/(99/100)<n$.

Stirling's bounds give $b(y_0)\le e^{1/(12n)}/\sqrt{2\pi(k-1)(1-y_0)}<47/1000$. The lattice sum of the unimodal kernel is at most its integral times $p$, plus its maximum. This follows by omitting the largest sample and charging each other sample to the adjacent interval toward the largest one. The beta integral is $1/(n+1)$. Since $p/(n+1)<1/200$, the total relative cost is at most
\[
\frac{61}{10}\left(\frac{47}{1000}+\frac1{200}\right)+\frac1{10{,}000}
<\frac8{25}.
\]
This proves (F4).

Equations (F1), (F3) and (F4) also give
\[
\frac{|\alpha_{n+1,p}-\alpha_{n,p}|}{\alpha_{n,p}}\le \frac{33}{100p}
\quad(1\le p\le J).
\tag{F5}
\]
Indeed the two boundary-error terms in (F1), multiplied by $p/\alpha_{n,p}$, cost at most $2 n^{3/2}e^{-(n-1)/\sqrt n}<1/100$ by (F2). The same upper bound for the two row sizes follows by using $\lfloor n/2\rfloor$ in both boundary terms. For $p=1$ the ratio between the two coefficients $\alpha$ is exactly one.

\subsection{Isolation at the adaptive radius}

Fix $1\le h\le H$. Put $\rho=\alpha_{n,h}/\alpha_{n,h+1}$, $u=\rho(h/(h+1))^n$ and
\[
F(z)=\frac{E_{n,k}(uz)}{\alpha_{n,h}h^n(uz)^h},
\qquad
b_p^*=\left(\frac ph\left(\frac h{h+1}\right)^{p-h}\right)^n.
\]
The terms at $p=h,h+1$ are exactly $1,z$. On either side of these terms the ratios of consecutive $b_p^*$ away from the central pair are at most $q=e^{-n/(h+1)^2}$. This follows by evaluating the largest ratios at the two immediate neighbors, where the exact bases are $1-1/h^2$ and $1-1/(h+1)^2$. The left tail is empty for $h=1$.

From (F2), for $p,v\le J$,
\[
\frac{\alpha_{n,p}}{\alpha_{n,v}}
\le \frac{51}{10}\max(1,v/p).
\]
Thus both $\rho$ and $\rho^{-1}$ are below 6 if $h\ge6$, and below $51/5$ for $h\le5$. In the first case $h+1\le(9/25)\sqrt n$, so $q<1/2200$. In the second case $q\le e^{-10000/36}$, which is more than enough for $q\max(\rho,\rho^{-1})<6/2200$ again. At $p=h-\ell$ the coefficient ratio to $\alpha_{n,h}$ is at most $(51/10)(\ell+1)$. At $p=h+1+\ell\le J$ the ratio to $\alpha_{n,h+1}$ is at most $51/10$. When $p=h+1+\ell>J$, we have $\ell>(7/50)\sqrt n$ and $h+1<(18/7)\ell$. Hence
\[
1/\alpha_{n,h+1}\le\frac{100}{99}(1+(8/3)(h+1))<8(\ell+1).
\]
These bounds, and $\alpha_{n,p}\le1$, give right-tail ratio at most $8(\ell+1)$ for every $p=h+1+\ell$.

Let $s$ be any upper bound on $q\max(|\rho z|,|\rho z|^{-1})$. At left index $h-\ell$ the absolute term is at most $(51/10)(\ell+1)s^\ell$, and at right index $h+1+\ell$ it is at most $8|z|(\ell+1)s^\ell$. Thus one can use $s=1/300$ on $|z+1|=1/8$, $s=3/500$ on $|z|=1/2$ and $s=9/1000$ on $|z|=3$. On these circles respectively the error in $F-(1+z)$ is less than $1/10$, $1/2$ and 2. These follow by multiplying $(51/10)+8|z|$ by $\sum_{\ell\ge1}(\ell+1)s^\ell=(1-s)^{-2}-1$. Rouché's theorem for $z^{h-1}F(z)$ gives respectively one root in the disk about $-1$, $h-1$ roots inside the half-radius circle and $h$ inside the radius-three circle. The disk sits in the resulting annulus. Its unique root is simple and negative by conjugation, and has absolute rank $h$.

At this root write $x=uz$, $X=|x|$. We have $|z+1|<1/8$. Since $F(z)=0$, the identity $xE_{n,k}'(x)/(\alpha_{n,h}h^n x^h)=zF'(z)-F(z)$ gives the two central terms $-1,0$. The remaining terms are bounded by
\[
\frac{51}{10}\left(\frac{1+s}{(1-s)^3}-1\right)
  +9\frac{2s}{(1-s)^3}<\frac3{20}\qquad(s=1/300).
\]
The two weights are $(\ell+1)^2$ on the left and $\ell(\ell+1)$ on the right, including the earlier coefficient-ratio factors. Thus $|xE'_{n,k}(x)|>(17/20)\alpha_{n,h}h^nX^h$.

\subsection{Signs and the join}

In the scale $\alpha_{n,h}h^n X^h$, (F3) and (F4) bound the contribution of indices $p\le J$ to $|xL(x)|$ by at most $(8/25)$ times the absolute sum of the terms of $F(z)$, which is less than $(8/25)(1+9/8+1/10)$. For $p>J$ use $\lambda_p\le p$, $\ell=p-h-1\ge14$, $p<4\ell$, and the bound $1/\alpha_{n,h+1}<8(\ell+1)$. Their normalized sum is less than $36\sum_{\ell\ge14}\ell(\ell+1)(1/300)^\ell<1/1000$. Therefore $|xL(x)|<(18/25)\alpha_{n,h}h^nX^h<|xE'_{n,k}(x)|$. At an old root $xP_n'(x)=E'_{n,k}(x)$, so $|L(x)|<|x(1-x)P_n'(x)|$. The latter term has sign $(-1)^h$ because all $h-1$ smaller roots are negative and simple, and $P_n(0)=1$. Both signed insertions preserve this sign.

For consecutive order, let $u_{n,h}$ and $u_{n+1,h}$ be the two adaptive radii. At $p=h,h+1$, (F5) bounds $\alpha_{n+1,p}/\alpha_{n,p}$ between $83/100$ and $117/100$ for $p\ge2$. At $p=1$ it equals one. The old root is therefore in the new rank-$h$ annulus $(u_{n+1,h}/2,3u_{n+1,h})$ in modulus. For $h\ge2$, its modulus divided by $u_{n+1,h}$ is between $(7/8)(83/117)$ and $(9/8)(117/83)(3/2)$. At $h=1$, it is between $(7/8)(83/100)2$ and $(9/8)(117/100)2$. The new polynomial has exactly one root in that annulus and sign $(-1)^{h-1}$ before that root. Its sign $(-1)^h$ at the old root places the old root strictly after it in modulus, with every other larger new root outside the annulus.

Finally every old front root has $X<(9/8)(51/5)(H/(H+1))^n$. At the next integer phase sample $\theta=(H+1)\pi/n$, one has $\pi\cot\theta<n/(H+1)$ and
\[
n\log(1+1/H)-\frac n{H+1}
\ge\frac n{2(H+1)^2}
\ge\frac1{2(9/25)^2}
>\log((9/8)(51/5)).
\]
Hence each of the $H$ front roots comes before this sample. The sample satisfies $n\sin^2\theta>6/5$ by $H+1>7\sqrt n/20$, $\sin y\ge y-y^3/6$ and $n\ge10{,}000$. Reflection gives the corresponding order and insertion signs at large modulus. The estimate leaves the middle-root argument as the remaining analytic step for this cutoff.

\section{Partial fractions for any number of occupied Eulerian poles}

This estimate compares the middle roots of consecutive rows with arbitrary numbers of permitted cycles. Let $S=\sin\theta$, $nS^2\ge1$, $n\ge100$, $n\ge k\ge32$. Put $D_0=\sqrt{n/k}$ and $U=3+\log(3D_0)$. Assume
\[
kS\ge2000 U. \tag{PA}
\]
We prove that the integer-phase signs and both insertion signs at old roots hold under (PA), within the residual-index range below. It suffices below to consider $n\ge2k-1$, since the earlier transitions are single-cycle Eulerian. For the lag estimates below first take $k\le n/3$; the other thresholds are covered by the one- and two-cycle band, apart from the single transition $n=3k-1$ which has $k>n/3$. The same residual estimates in fact hold for $k\le(n+3)/3$, so the displayed proof covers this transition too.

\subsection{Expanding by occupied pole subsets}

Use $z_p=(\cos\theta+ipS)^{-1}$, $\rho_p=|z_p|$, and $F_j,D,d_j$ as in the \href{https://github.com/ProgramComputer/open-math-research/blob/main/topics/long-cycle-excedances/central-interval.md}{central interval proof}. Put $g_j=d_j$ for $j\ge k$ and zero elsewhere. Let $b_j=g_j-g_{j-1}$, so $\sum|b_j|<2/k$ and $b_j=0$ below $k$. The exact generating series is
\[
\sum_m F_m u^m=\prod_{p\,\mathrm{odd}}D(z_pu),\qquad
D(z)-1=\frac{\sum_{j\ge k} b_j z^j}{1-z}.
\]
The contribution supported on a single pole is $d_m B_m$, $B_m=\sum_p z_p^m$. Write $R_m=F_m-d_m B_m$ when $m\ge k$. For a finite set $I$ of $h\ge2$ distinct occupied poles the coefficient with total exponent $m$ is summed over $s_p\ge k$, $\sum s_p\le m$, with weight $\prod b_{s_p}$. By partial fractions of $\prod_{p\in I}(1-z_pu)^{-1}$ its remaining factor is
\[
\sum_{a\in I} \prod_{p\in I} z_p^{e_p}
\prod_{q\in I\setminus\{a\}}(1-z_q/z_a)^{-1},
\]
where $e_q=s_q$ for $q\ne a$, $e_a=s_a+m-\sum s_p$, and $\sum e_p=m$. The distinctness of the $z_p$ holds for $0<\theta<\pi$. In particular $e_p\ge k$ and some $b\in I$ has $e_b\ge m/h$.

We will bound the value for $m=n$ and $m=n+1-k\ge k$ with $mS^2\ge2/3$. The derivative bound only uses $m=n$. Write $w_p=\rho_p^{k/3-1}$. Since
\[
|1-z_q/z_a|=|q-a| S\rho_q,
\]
the absolute value of each term inside the sum over $a$, after suppressing the $b_{s_p}$ weights, is at most
\[
S^{1-h}\rho_b^{m/(3h)}\,w_a\prod_{q\ne a}\frac{w_q}{|q-a|}. \tag{PB}
\]
We took $k/3$ powers at every pole and $m/(3h)$ at $b$. A further $e_p/3$ powers remain at each pole and can be used for differentiating. We can upper bound (PB) by summing over all choices of $b$ before summing the weights. Their absolute sum for a fixed $I$ is at most $(2/k)^h$.

For the normalized derivative at $m=n$, differentiating the powers costs at most $9h/4$ in this bound. For $|p|=1$, the cost $e_p/n$ is at most one. For $|p|\ge3$, $|z_p'/z_p|\le|p|\rho_p$ and
\[
\frac{e_p}{n}|p|\rho_p^{1+e_p/3}
\le \frac{6|p|}{n(p^2-1)S^2}\le\frac94
\]
by Bernoulli's inequality with exponent $e_p/6\ge1$. Each differentiated inverse denominator has relative magnitude $\rho_q/S$, since it equals $(\cot\theta+iq)/(i(q-a))$. Thus differentiation of the denominators costs at most $(h-1)/(nS)$, giving total normalized derivative cost below $5h/2$. These bounds apply to the finite-degree coefficient for each finite set and pass to the full pole sum by local absolute convergence.

\subsection{Uniform summation of the partial fractions}

For $\lambda\ge 8$ use $(1+x)^{\lambda/2}\ge(1+(\lambda/4)x)^2$. For $|p|\ge3$, $p^2-1\ge8p^2/9$. Monotone integral comparison for the positive odd integers, then reflection, gives
\[
\sum_p\rho_p^\lambda\le2+\frac{3\pi}{4\sqrt2}\frac1{S\sqrt\lambda}.
\]
Therefore $\sum_p w_p\le2+4D_0$, and $V_h=\sum_p\rho_p^{m/(3h)}\le 2+4\sqrt h\le6h$ for $mS^2\ge2/3$ and $h\le m/k$. For each odd $a$ split the sum of $w_q/|q-a|$ at distance $2\lceil2D_0\rceil$. The near terms are bounded by the harmonic sum through $\lceil2D_0\rceil$, and the far ones by $(2+4D_0)/(4D_0)$. Hence $\sum_{q\ne a}w_q/|q-a|\le U$.

In summing (PB), the part $a=b$ costs at most $V_h U^{h-1}/(h-1)!$ after omitting $S^{1-h}$. For $b\ne a$, first sum the $h-2$ remaining poles for fixed $a,b$. The sum of $w_a/|a-b|$ over $a\ne b$ is at most $U$ for each $b$, and $\sum_b\rho_b^{m/(3h)}w_b\le V_h$. This gives at most $V_h U^{h-1}/(h-2)!$. Keeping the largest exponent at some pole is important here, since otherwise $V_h$ would depend on $n/k$. The two bounds together are $h V_h U^{h-1}/(h-1)!$.

Put $\eta=2U/(kS)\le1/1000$. Since $h^2\le 4^{h-1}$ and $h^3\le8^{h-1}$,
\[
|R_m|\le \frac{12}{k}\frac{4\eta}{1-4\eta}<\frac1{20k},\qquad
\frac{|R'_n|}{n}\le \frac{30}{k}\frac{8\eta}{1-8\eta}<\frac1{4k}. \tag{PC}
\]
All remaining poles and every number of cycles are included.

\subsection{Applying the old-root insertion comparison}

For $n\ge100$, $nS^2\ge1$ the ordinary single-pole tail obeys $\sum_{p=3,5,\ldots}p\rho_p^n<1/10$, by equation (6) of the \href{https://github.com/ProgramComputer/open-math-research/blob/main/topics/long-cycle-excedances/first-lag.md}{first two-cycle boundary argument}. Thus $|B_n-2\cos(n\theta)|<1/15$ and $|B'_n/n+2\sin(n\theta)|<1/5$. With $1/(2k-1)\le d_n\le1/k$ and (PC), the total value error is below $(1/15+1/20)/k$ and the derivative error below $(1/5+1/4)/k$. This gives the integer samples and, at every old zero, $|\cos(n\theta)|<7/60$, $|\sin(n\theta)|>19/20$, and $|F'_n|>n/(2k)$.

At a zero of $F_n$ the normalized insertion contributions are $S F'_n$ and $B_k F_m$ where $m=n+1-k$. If $mS^2\ge2/3$ and $m\ge k$, the absolute sum bound above at $\lambda=m$ yields $|B_m|\le2+3\pi\sqrt3/8< 4.1$. Thus $|F_m|<5/k$. Also
\[
\frac{|B_k|}{nS}\le \frac2{nS}+\frac{3\pi}{4\sqrt{2k}\,nS^2}
\le \frac{2}{2000}+\frac{3\pi}{4\sqrt{4000}}<\frac1{10},
\]
using $n\ge k$, $kS\ge2000$ and $nS^2\ge1$. Consequently $|B_kF_m|<nS/(2k)<S|F'_n|$.

Taking $k\ge2000\sqrt n\,[3+\log(3\sqrt{n/k})]$ suffices for (PA) throughout $nS^2\ge1$. This is an explicit effective but conservative condition. It leaves the intermediate threshold regime open.

\section{An explicit strict-interlacing range with unbounded cycle count}

For fixed $3\le k\le n$, let $E_{n,k}(t)$ count permutations of $[n]$ by excedances subject to every cycle having at least $k$ letters. Write $P_n=E_{n,k}/t$. Inversion of a permutation exchanges excedances and deficiencies. A single cycle can attain either extreme, so $P_n$ is reciprocal of degree $n-2$ with constant and leading coefficients one. \href{https://www.symmetricfunctions.com/realRootedOpenProblems.htm\#conjLongCycleExcedances}{Original conjecture} $\cdot$ \href{https://github.com/ProgramComputer/open-math-research/blob/main/topics/long-cycle-excedances/two-cycle-band.md}{Cycle weights and recurrence}

\textbf{Partial theorem.} The zeros of $P_n$ and $P_{n+1}$ are simple, negative and strictly interlace whenever $n\ge10{,}000$ and
\[
k\ge2000\sqrt n\left(3+\log\left(3\sqrt{n/k}\right)\right).
\]
The bound is effective and permits unboundedly many cycles. For example, take $n=2^{4j}$ and $k=2^{3j}$ for any integer $j\ge15$. The condition reduces to $2^j\ge2000(3+\log3+(j/2)\log2)$. Already $2^{15}>2000(5+15/2)$, and doubling the left side at each step preserves the inequality. The largest allowed cycle count is then $2^j$.

\subsection{Counting every root}

If $n\le3k-2$ this is the \href{https://github.com/ProgramComputer/open-math-research/blob/main/topics/long-cycle-excedances/two-cycle-band.md}{one- and two-cycle theorem}. In the remaining range $n\ge3k-1$, the residual index $m=n+1-k$ is at least $2n/3$ and $m\ge k$. Put $H=\lfloor7\sqrt n/20\rfloor$. The \href{https://github.com/ProgramComputer/open-math-research/blob/main/topics/long-cycle-excedances/effective-front.md}{coefficient argument} locates the $H$ smallest roots in modulus, all negative and simple, and verifies both insertion signs there. The last of these roots in modulus precedes the sample of index $H+1$.

Parameterize the negative axis by $t=-\exp(-\pi\cot\theta)$ with $0<\theta<\pi$. Then $|t|$ increases with $\theta$. On the first half of the interval take samples $\theta_j=j\pi/n$, starting with $j=H+1$. At this first sample $n\sin^2\theta_j>6/5$. All angles through $\pi/2$ consequently have $n\sin^2\theta\ge1$, $m\sin^2\theta\ge2/3$ and $k\sin\theta\ge2000(3+\log(3\sqrt{n/k}))$. The \href{https://github.com/ProgramComputer/open-math-research/blob/main/topics/long-cycle-excedances/separated-poles.md}{separated-pole estimate} gives the sample sign $(-1)^j$ for $E_{n,k}$ and controls the insertion term at every old zero in the interval. Reciprocity extends the sample signs, the old-root comparison, and the exterior-root count to the other half. For odd $n$ the central sample gap straddles $\pi/2$ and contains the reciprocal root $-1$.

The consecutive samples with indices $H+1,\dots,n-H-1$ give $n-2H-2$ sign changes. The $2H$ distinct exterior roots exhaust the remaining degree. Thus every gap contains precisely one root and every root is simple. At every old root $x$ we have
\[
P_{n+1}(x)=x(1-x)P'_n(x)+L(x),\qquad
|L(x)|<|x(1-x)P'_n(x)|.
\]
The signs of $P_{n+1}$ there are successively $-1,+1,\dots$ in order of increasing modulus. The positive constant term and the positive leading coefficient, together with $\deg P_{n+1}=n-1$, give both outside roots as well as one new root in each old gap. This proves the strict interlacing.

\section{A numerical cutoff for every minimum cycle length}

\textbf{Partial theorem.} For $n\ge10^{20}$ and every $3\le k\le n$, the nonzero zeros of the two consecutive minimum-cycle-length excedance polynomials are simple, negative and strictly interlace. This makes the earlier existential cutoff explicit. Checking the entire finite complement by enumeration is computationally impractical. The sign comparison below is for the actual next row.

Use the \href{https://github.com/ProgramComputer/open-math-research/blob/main/topics/long-cycle-excedances/effective-threshold.md}{effective large-threshold theorem} if $k\ge2000\sqrt n(3+\log(3\sqrt{n/k}))$. In the remaining range the short-cycle estimate applies. More generally, it proves strict interlacing for every $n\ge10^6$ satisfying
\[
\frac n{200k}>30\log n+20. \tag{U1}
\]
To check that the two sufficient conditions cover every threshold when $n\ge10^{20}$, observe first that if $k\le\sqrt n$ equation (U1) follows directly. Otherwise $k>\sqrt n$ and $k<2000\sqrt n(3+\log 3+(\log n)/4)$. At $n=10^{20}$, $400000(3+\log3+(\log n)/4)(30\log n+20)<10^{10}$ using $\log 3<6/5$ and $\log10<7/3$. The ratio of $\sqrt n$ to these two logarithmic factors increases for $\log n\ge40$, by logarithmic differentiation.

\subsection{A common coefficient contour}

Use $t=-e^{-s}$, $s\ge0$, $w_p=s+ip\pi$ for odd $p$, $r=|w_1|$, $S=\sin\theta=\pi/r$ and $z_p=r/w_p$. Define $T_s$ by the \href{https://github.com/ProgramComputer/open-math-research/blob/main/topics/long-cycle-excedances/fixed-threshold.md}{fixed-threshold coefficient identity}
\[
\frac{E_{j,k}(t)}{j!(1-t)^{j+1}}=[w^j]\frac{T_s(w)}{1+e^{w-s}}.
\]
The logarithm $G_s(w)$ of $T_s(w)$ is the specific polynomial with $G_s(0)=0$. Its linear term is $e^{-s}w/(1+e^{-s})$. Its terms of degree $\ell\ge2$ have absolute size at most $(|w|/r)^\ell\sum_p|z_p|^\ell/\ell$, by the Eulerian pole formula. For $s>0$ it also has the absolutely convergent expansion
\[
G_s(w)=-\sum_{a\ge1}\frac{(-e^{-s})^a}{a}\sum_{\ell=1}^{k-1}\frac{(aw)^\ell}{\ell!}.
\]
On $|w|\le R_*=r(1+1/(50k))$ we claim
\[
|G_s(w)|\le B:=6+10\log k . \tag{U2}
\]
For $s\le10$, $\sum_p|z_p|^2\le2+1/(2S^2)<8$, by telescoping $\sum_{p=3,5,\dots}1/(p^2-1)=1/4$. The linear term costs less than six, and $8e^{1/50}<10$ bounds the coefficient of the harmonic sum in the rest.

For $s>10$ put $L=\min(k-1,\lfloor s^2/16\rfloor)$ and $\delta=R_*-s\le5/s+s/(40k)$. For terms of degree at most $L$, use the exponential sum. The absolute truncated exponential, including the factor $e^{-as}$, is bounded by $e^{a\delta}\Pr(\operatorname{Pois}(aR_*)\le L)$. If $a\le2L/s$, then $a\delta\le5/8+1/20<1$, so this range contributes at most $3(1+\log k)$. For $a>2L/s$ the truncated exponential series times $e^{-aR_*}$ is bounded by the lower-tail probability $\Pr(\operatorname{Pois}(aR_*)\le aR_*/2)\le e^{-aR_*/8}$. Since $\delta/s<3/50$, the resulting sum is at most $\sum_{a\ge1}e^{-as/16}/a<2$. Any remaining degrees have $\ell\ge7$ and $\ell S^2>1/2$. By the lattice integral estimate in \href{https://github.com/ProgramComputer/open-math-research/blob/main/topics/long-cycle-excedances/separated-poles.md}{the pole argument},
\[
\sum_p|z_p|^\ell\le2+\frac{3\pi}{4\sqrt2 S\sqrt\ell}<5 .
\]
Here the integral estimate is valid for $\ell\ge4$ because $(1+x)^{\ell/4}\ge1+\ell x/4$. Thus the higher degrees cost at most $6\log k$, proving (U2).

Take $nS^2\ge1$. If $s\le10\sqrt k$ use the circle $R=r(1+1/(100k))$. Only $w_{\pm1}$ are inside, since $R_*^2-r^2<5$ whereas $|w_3|^2-r^2=8\pi^2$. The distance from the circle to each pole exceeds $n^{-2}$, using $r/(100k)\ge\pi/(100n)$ for the first pair and the squared gap for the others. If $s>10\sqrt k$, choose a radius in $[r(1+1/(100k)),R_*]$ whose distance from every pole radius is at least $n^{-2}$. This interval has length at least $\pi/(100n)$ and fewer than $10\sqrt n$ pole radii can meet it within $n^{-2}$, as $r\le\pi\sqrt n$. On both circles $|1+e^{w-s}|^{-1}\le4n^2$. Indeed, when $|\Re(w-s)|\ge n^{-2}/2$ use $|1-e^{\Re(w-s)}|$. Otherwise the distance of $\Im w$ from the nearest odd multiple of $\pi$ is at least $\sqrt3\,n^{-2}/2$, and the squared modulus is $(1-e^{\Re(w-s)})^2+4e^{\Re(w-s)}\cos^2(\Im w/2)$.

For $j=n,n+1$ use the common positive normalization
\[
Y_j=\frac{r^{j+1}E_{j,k}(t)}{2|T_s(w_1)|j!(1-t)^{j+1}} .
\]
The residues give the sum over enclosed positive odd $p$ of $\Re(T_s(w_p)z_p^{j+1})/|T_s(w_1)|$. The circle remainder in either row has magnitude at most $2n^2r\exp(2B-n/(200k))<n^{-2}$, by (U1).

\subsection{Neighboring residues}

Put $\gamma=\Im G_s(w_1)$ and $\Psi=(n+1)\theta-\gamma$. The principal pair contributes $\cos\Psi$ to $Y_n$. For $s>10\sqrt k$, at odd poles one has
\[
G_s(w_p)=-\log(1+e^{-s})-H(w_p),\qquad
H(w)=\sum_{a\ge1}\frac{e^{-aw}}a\sum_{\ell=0}^{k-1}\frac{(aw)^\ell}{\ell!}.
\]
This follows by substituting $e^{-aw_p}=(-e^{-s})^a$ in the polynomial expansion. On the vertical segments from $s$ to enclosed $w_p$ one has $\log(|w|/s)\le 7/(100k)$. Differentiating the uniformly converging series for $H$ there gives
\[
|H'(w)|\le e^{7/100}\sum_{a\ge1}e^{-as}(as)^{k-1}/(k-1)!.
\]
For $k\ge14$ the sum is at most $1/s+1/\sqrt{2\pi(k-1)}$, by summing the unimodal gamma density on a mesh of spacing $s$. Its integral is one and Stirling's lower bound gives the stated bound at its mode $k-1$. Thus $\pi |H'|\le\pi e^{7/100}(1/(10\sqrt{14})+1/\sqrt{26\pi})<12/25$. For $3\le k\le13$, $s>10\sqrt k\ge2k$ and $s>17$. Each summand is bounded by the Poisson probability of falling below half its mean, at most $e^{-as/8}$. Therefore $\pi |H'|\le\pi e^{7/100}/(e^{17/8}-1)<12/25$. Since $H(s)$ is real,
\[
|\gamma|<12/25,\qquad
|G_s(w_p)-G_s(w_1)|\le(12/25)(p+1) \quad(p\ge3).
\]

An enclosed pole has $(p^2-1)S^2\le1/(20k)$, hence $|z_p|^n\le\exp(-(49/100)(p^2-1))$. The total nonprincipal weighted error on this circle is bounded by
\[
\sum_{p=3,5,\ldots}p\,|z_p|^n\,|T_s(w_p)/T_s(w_1)|
<3e^{-2}+\sum_{p=5,7,\ldots}p e^{-(37/100)(p^2-1)}<21/50 ,
\]
where the sum on the left runs only over enclosed poles. For $p\ge5$, $(p+1)/(p^2-1)\le1/4$. The remaining sum is bounded by its first term and half the integral of the decreasing summand from $5$ to infinity. For $s\le10\sqrt k$ the nonprincipal sum is empty. Subtracting $\cos\theta\,Y_n$ from the following row multiplies each pole contribution by $z_p-\cos\theta$, whose magnitude is at most $pS$. Thus
\[
|Y_n-\cos\Psi|<1/2,\qquad |Y_{n+1}-\cos\theta\,Y_n+S\sin\Psi|< S/2. \tag{U3}
\]
These hold throughout $nS^2\ge1$. The correction $\gamma$ need only be bounded on the outer endpoint for the root count.

\subsection{A phase-crossing argument}

Let $H_0=\lfloor7\sqrt n/20\rfloor$ and $\theta_*=(H_0+5/6)\pi/n$. Then $n\sin^2\theta_*>1$, $\theta_*<1/100$ and $s_*=\pi\cot\theta_*>2\sqrt n>10\sqrt k$. These follow from $H_0+1\le(9/25)\sqrt n$, the sine and cosine Taylor bounds, and (U1). Thus $\Psi(\theta_*)<(H_0+1)\pi$ by $12/25+1/100<\pi/6$. All $H_0$ exterior roots from the \href{https://github.com/ProgramComputer/open-math-research/blob/main/topics/long-cycle-excedances/effective-front.md}{effective front proof} lie below $|t(\theta_*)|$. Indeed their moduli are bounded by $(9/8)(51/5)(H_0/(H_0+1))^n$ and
\[
n\log(1+1/H_0)-\frac n{H_0+5/6}
\ge \frac n{(H_0+1)^2}\left(\frac12-\frac{H_0+1}{6(H_0+5/6)}\right)
>\frac{33/100}{(9/25)^2}>\frac52>\log((9/8)(51/5)).
\]
Use $\pi\cot\theta_*<n/(H_0+5/6)$. At the central endpoint $s=0$, $\gamma=\pi/2$: the linear term gives this value, even powers in $G_0(i\pi)$ are real, and odd powers of degree at least three have zero coefficient by the Eulerian pole formula. Hence $\Psi(\pi/2)=n\pi/2$ exactly.

Take the first hits, after $\theta_*$, of the levels $j\pi$ for $H_0+1\le j\le\lfloor n/2\rfloor$. They occur in the indicated order by continuity. Between two consecutive hits take the last hit of the lower level before the upper one. On the intervening open interval the phase is strictly between the two levels. By (U3) the old polynomial has opposite endpoint signs. Select one zero there. Its phase gives the sign of the next polynomial by (U3), since $|\cos\Psi|<1/2$ implies $|\sin\Psi|>\sqrt3/2$. For even $n$, an early first hit of the central level is harmless: the reciprocal intervals are disjoint and the degree count below exhausts all roots. For odd $n$, the odd reciprocal polynomial $P_n$ has the additional zero $-1$ with the required next-row sign directly from (U3) at $\pi/2$. Reflecting the selected half-intervals and including the $2H_0$ exterior roots yields $n-2$ distinct negative roots for $P_n$ in both parities. Its degree is $n-2$, so each is simple. The old-root signs from (U3) and the exterior-root estimates alternate for $P_{n+1}$ after division by $t$. Reciprocity gives $P_{n+1}(1/x)=x^{-(n-1)}P_{n+1}(x)$ at the reflected roots. Its constant and leading coefficients are one and its degree is $n-1$. The intermediate value theorem gives strict interlacing. This proof needs no derivative bound or monotonicity for the phase correction.


\section*{Reproducibility and scope}
The complete dyadic certificate and three independently usable checkers are in the
\href{https://github.com/ProgramComputer/open-math-research/tree/research/topics/long-cycle-excedances}{topic archive}.
A second implementation recovered all row coefficients by inclusion--exclusion. An independent agent-assisted referee
recomputed all 79,600 endpoint signs by clearing denominators and checked the analytic argument.
The proof presented here is on paper. No end-to-end Lean or Rocq verification is asserted.
\end{document}
